\chapter{The Square-Root Door and the Universe Kernel}
\label{chap:sqrt-door-universe-kernel}

The operator of the last chapter measured with an energy, and an energy is a squared quantity --- it
returns the square of a length, not the length itself. This chapter takes the square root. Doing so
turns the energy into a genuine norm, a real measure of size, and in the same motion arrives at a
threshold the construction has been approaching since its second chapter: the door to a completed
Hilbert space, the infinite-dimensional home in which the construction's approximations would finally
converge. The chapter is honest about that door. It walks up to it, takes the square root that makes a
norm, and marks plainly the point where the completion to a real Hilbert space would be imported --- it
opens the door and declines to claim it has built the room. Then it does one more thing the rest of the
book depends on: it binds the single invariant to the labels physics gives matter, charge and color
and flavor, in an object it calls the universe kernel.

\section{The square-root door}
\label{se:the-square-root-door}

A norm is a length, and a length is the square root of an energy. The construction has, in the energy
of the previous chapters, a quantity that behaves like a squared length: it is non-negative, it adds
across independent contributions, and it detects zero. To get from it to an honest measure of size, the
construction introduces a square root and defines the norm of a configuration as the square root of its
energy. The one property it requires of that square root is that it be zero-detecting --- that the
square root of a quantity is zero only when the quantity is --- and from that the norm inherits exactly
what it needs: the norm of a configuration is zero precisely when the configuration is the zero
configuration. The energy detected zero; the square root passes the detection through; the norm detects
zero too. The construction has its first genuine length.

It is worth pausing on what a length buys that an energy did not. An energy compares a configuration
only with itself --- it returns one number, the configuration's own size squared. A norm, being a
length, lets the construction compare two configurations: the distance between them is the norm of their
difference, and with distances in hand the configuration space becomes a metric space, a place where
nearness and convergence have exact meanings. This is the step that makes the variational approximations
of Part II more than a sequence of numbers. Once there is a distance, a sequence of finite solutions
that agree more and more closely is a Cauchy sequence in a precise sense, and the question of whether it
converges becomes the question of whether its limit point is present in the space. The norm turns the
loose talk of approximations getting closer into the exact machinery of metric convergence, and it is
the last thing the construction needs before it can ask, rigorously, where its limits live.

The square root the construction uses is itself worth a word, because the construction is careful not to
assume more of it than it needs. It does not build the real square-root function --- that would require
the very completion it is trying to defer. It assumes only an abstract square-root operation with the one
property that matters, that it sends zero to zero and nothing else to zero, and it notes that the
plainest instance of such an operation is the trivial one over the whole numbers, where the square root
is left undone and the energy stands in for its own length. Between that degenerate shadow and the real
square root over a completed field lies exactly the gap the chapter is honest about: the construction can
always take the trivial square root and stay finite, and it can name the real one and reason with it, but
the real one lives on the far side of the door.

Taking the square root is also stepping through a door, and the construction names the door rather than
pretending it is not there. A norm makes a configuration space into a metric space, a space with
distances, and a metric space has a completion --- the larger space got by filling in the limits of all
its Cauchy sequences. For the construction's norm that completion is a Hilbert space, the
infinite-dimensional space of square-summable configurations in which the variational approximations of
Part II would have their limits. But the completion is exactly the move a finite construction cannot
finish: filling in all the limits is filling in a continuum's worth of points the construction can only
ever approach. So it does the honest thing. It packages the norm together with an explicit token --- a
marker that says, \emph{here is where the completion to a Hilbert space is imported} --- and it imports
nothing it has not built. The construction goes exactly as far as a finite procedure can, to the edge
of the metric space and the threshold of its completion, and it flags the threshold instead of
stepping across it. This is the same honesty the number tower showed in the second chapter, where a real
was a sample and the continuum was approached but never completed; the Hilbert space is the continuum
again, one dimension up, and the square-root door is where the construction marks how close it has come
and where the rest would have to be taken on the strength of a completion it gestures at but does not
construct. The experiment on Hilbert's effect states the destination precisely: the space of admissible
spline refinements, completed under prediction and consistency, \emph{is} a Hilbert space, its inner
product induced by informational minimality. The construction agrees, and adds only the finitist's
footnote --- that completing is a thing done at the door, in the limit, up to any epsilon one demands
and never at a finite stage.

The completion itself is worth describing even though the construction does not perform it, because
naming what is deferred is half the honesty. To complete a metric space is to funge into it all the
limit points its Cauchy sequences are converging toward --- to pool the space together with the
destinations of its convergent sequences, so that every sequence that ought to converge has somewhere to
converge to. For a finite-dimensional space this adds nothing; for the construction's infinitely-
refinable configuration space it adds a continuum of new points, the square-summable configurations no
finite refinement ever reaches but that all the refinements approach. That pooled limit space is the
Hilbert space, and it is genuinely larger than anything the construction builds. The construction's
discipline is to use the Hilbert space's name and properties --- to reason as though the limits were
there --- while marking, at the door, that they were imported on the strength of a completion it
gestured at, not assembled from Facts. It borrows the room and pays for it with a clearly posted note.

\section{The universe kernel}
\label{se:the-universe-kernel}

The second move of the chapter is the one that turns the invariant toward physics, and it is worth
watching closely because it is the first time the single invariant of Part II acquires the labels of
matter. The construction builds an object it calls the universe kernel: the same invariant, the same
positive-definite energy that detects zero, now carrying the quantum labels that physics hangs on a
particle --- a charge, a color, a flavor, a phase. The invariant does not change; what changes is that
it is now indexed, tanged by these characteristics --- some of them read off as observable values, others
only carried --- so that the one number can be tracked for each characteristic it bears. And the properties earned over the last several chapters carry to the top
unbroken: the universe kernel still detects zero, reporting nothing exactly when nothing is present, and
it is still coercive, its energy still controlling the size of what it measures, so that the
well-posedness established on three nodes holds for the labeled invariant as well.

The indexing is more concrete, and more honest, than the grand names suggest --- and it is not the same
operation for every label. To bind the invariant to a \emph{charge} is to let the single invariant be
read off for that charge the way one ledger tracks a currency: the electric charge, and the phase that
travels with it, are \emph{observable} --- read out of the kernel as definite values, the sign and the
clock a measurement can actually report. Color and flavor are not like that. They are \emph{carried} ---
borne by the kernel as characteristics it genuinely possesses and conserves --- but they are never read
out on their own; the construction attaches a color and a flavor and preserves them without ever exposing
a bare color or a bare flavor to a reading. Nothing about the energy changes; it is the same
positive-definite, zero-detecting form. What changes is that the construction can now ask how much
invariant is present \emph{of a given charge} and get a definite number, while color and flavor remain
present, conserved, and unread. That is a sharper claim than a label drawn from a tiny finite set, and a
more honest one: it is the shape of confinement --- a characteristic the apparatus carries and
structurally cannot read out. That charge and phase are observable while color and flavor are only
carried is the modeling claim; that an invariant can be cleanly indexed by conserved characteristics,
some read off and some merely borne, is the mathematics.

Binding the invariant to charge is what lets a force fall out of it, and the experiment on the
inverse-square law shows how. The influence of a refinement event, that experiment finds, decreases as
the inverse square of the informational separation between events --- and this is not postulated but
forced, by the geometry of admissible splines and the law of spline sufficiency that governs them. Read
through the universe kernel, the inverse-square falloff is the shape the single invariant takes when it
is carried by a charge across a distance: the coercive, zero-detecting energy, spread over the spline
geometry, produces exactly the inverse-square influence that charge is known to exert. The construction
is not deriving electromagnetism in full; it is showing that once the invariant is bound to a charge,
the falloff of that charge's influence is not a separate law to be added but a consequence of the
spline geometry already in hand. The source of such influence the construction marks with a name of its
own, a white-hole source, the place a charge's effect originates in the record.

It is worth being clear about how strong, and how limited, this is. The construction does not write down
Maxwell's equations and read the inverse-square off them; it shows that the influence one charged event
exerts on another, computed through the spline geometry the law of spline sufficiency forces, falls off
as the inverse square of their separation, with no separate assumption invoked to make it do so. That is
a real result --- the falloff is a theorem about the geometry, not a postulate --- and it is the first
time a recognizable physical force-law has emerged from the construction rather than been placed into
it. But it is also narrow: it is the falloff, the shape of the influence with distance, and not the full
dynamical theory of the field. The construction has shown the geometry already contains the
inverse-square, which is a great deal more than nothing and a great deal less than a derivation of
electromagnetism, and the honest reading sits exactly between those two.

One phrase from the construction's own notes deserves unpacking: coercivity at the top. The worry, when
an invariant is dressed in physical labels and carried to the most elaborate level the construction
reaches, is that the careful properties earned on three nodes might quietly fail somewhere along the way
--- that the labeled, large-scale invariant might develop a flat direction, a place where it stops
controlling the size of what it measures, and become ill-posed exactly where the physics gets
interesting. The construction proves this does not happen. Coercivity, established at the bottom on the
three-node form, is shown to survive all the way up to the universe kernel: the labeled invariant at the
top is still bounded below, still controls the size of its configurations, still well-posed. That the
good behavior reaches the top unbroken is what lets the later chapters read physical content off the
invariant without fearing that the reading is an artifact of a degeneracy hiding in the apparatus.

\section{What is demonstrated, and what is assumed}
\label{se:demonstrated-assumed-ch16}

The separation here is more interesting than usual, because this is the chapter where the construction
deliberately marks a boundary it does not cross.

What is demonstrated is the square-root norm and the labeled invariant. Given a zero-detecting square
root, the norm built as the square root of the energy is proved to detect zero, a genuine length. The
universe kernel is built by binding the invariant to charge, color, and flavor, and is proved to retain
both of the energy's load-bearing properties: it detects zero, and it is coercive. These are theorems
about finite objects, and they hold whatever names the labels are given.

\begin{quote}
\emph{A note on the labels --- and on the door.} Two things here are bridges and one is a marked gap.
The physical labels --- \emph{charge}, \emph{color}, \emph{flavor} --- are names hung on indices of the
invariant; the mathematics is the indexing and the preserved definiteness, and the identification of
those indices with the quantum numbers of matter is the modeling, earned to the degree the analogy
holds. The inverse-square reading of the invariant's falloff is likewise a bridge from the spline
geometry to a physical force. But the Hilbert completion is neither a theorem nor a modeling bridge ---
it is a gap the construction marks explicitly and does not fill. The completed Hilbert space is imported
at the door, not built; the construction proves everything up to the threshold and flags the threshold,
and a reader is owed the plain statement that the infinite-dimensional home of these approximations is
approached, named, and left uncompleted, exactly as the continuum has been throughout.
\end{quote}

What is assumed, then, is the completion at the door --- the one place this chapter reaches past what it
constructs --- together with the standing oracle and the physical labels. With a real norm in hand, a
labeled invariant that detects zero and stays coercive, and an honest marker at the threshold of the
Hilbert space, the construction has assembled the operator side of its account. The next chapters turn
from the operator to the record it writes: how a concrete operator grounds the whole abstract apparatus,
when an invariant is physical and when it is only an artifact of description, and what it means to write
down the first genuine measurement in the construction's own vocabulary.
